\section{Additional Downstream and Benchmark Details}

\subsection{Calibration remains secondary}

Table~\ref{tab:appendix_calibration} summarizes the main calibration readout
for the strongest visible and masked rerankers on \textsc{Craft-Deploy}. The
gap is much smaller than the ranking and exploitability gap in the main text,
which is why calibration remains an
appendix-level result rather than a core mechanism claim.

\begin{table}[h]
\centering
\caption{Calibration summary for the main deployment comparison. Lower ECE is
better.}
\label{tab:appendix_calibration}
\small
\begin{tabular}{lcc}
\toprule
model & ECE & note \\
\midrule
\texttt{reranker8\_masked} & 0.4195 & best deployment model overall \\
\texttt{reranker8\_visible} & 0.4320 & utility and faithfulness both weaker \\
\bottomrule
\end{tabular}
\end{table}

\subsection{\textsc{Craft-Clean} acceptance summary}

Table~\ref{tab:appendix_acceptance} records the acceptance logic for the
current \textsc{Craft-Clean} export. We keep the universal strict gate as a
diagnostic, but the benchmark is accepted under the artifact-focused criterion
that does not treat purely semantic detectability as a benchmark failure.

\begin{table}[h]
\centering
\caption{\textsc{Craft-Clean} acceptance summary. Higher accepted-family counts
are better; the strict criterion is reported only as a diagnostic.}
\label{tab:appendix_acceptance}
\small
\begin{tabular}{lcc}
\toprule
criterion & accepted families & interpretation \\
\midrule
strict & 2 / 54 & over-strict diagnostic only \\
\texttt{ignore\_semantic\_only} & 52 / 54 & authoritative export criterion \\
\texttt{surface\_or\_mixed\_only} & 52 / 54 & artifact-focused cross-check \\
\bottomrule
\end{tabular}
\end{table}

\subsection{\textsc{Craft-Clean} reproduction summary}

Table~\ref{tab:appendix_benchv3_repro} gives the cleaner-benchmark frozen-head
means across three seeds. These numbers support the same qualitative story as
the main text: \textsc{Craft-Clean} preserves process-faithfulness gaps, but it does
not amplify them into a strong utility separation.

\begin{table}[h]
\centering
\caption{\textsc{Craft-Clean} quartet-protocol reproduction summary across
three seeds. Higher is better for AUROC and AMCD; lower is better for ASS.}
\label{tab:appendix_benchv3_repro}
\small
\begin{tabular}{lccc}
\toprule
model & AUROC & AMCD & ASS \\
\midrule
\texttt{step\_only\_bce} & 0.9233 & 0.8842 & 0.0891 \\
\texttt{visible\_bce} & 0.9161 & 0.9028 & 0.1389 \\
\texttt{masked\_bce} & 0.9302 & 0.9213 & 0.1055 \\
\texttt{pairwise\_visible} & 0.8659 & 0.9028 & 0.0736 \\
\bottomrule
\end{tabular}
\end{table}

\subsection{\textsc{ProcessBench} external-source summary}

Table~\ref{tab:appendix_processbench} records the current external-source main
table behind the brief readout in the main text. The important point is not
that rerankers are useless; it is that the current stronger reranker does not
yet dominate the best frozen external-source baseline, and that its weakest
behavior on this benchmark is calibration rather than total lack of ranking
signal.
The package is now materialized in four canonical release artifacts:
\texttt{processbench\_suite.json}, \texttt{processbench\_suite.md},
\texttt{processbench\_manifest.json}, and
\texttt{processbench\_split\_manifest.json}. Together they fix the
authoritative datasets and result artifacts for \textbf{PB-Trace},
\textbf{PB-Prefix}, answer-only swap, and local repair. They also pin the
grouped trace split and the capped prefix split used by the current frozen and
reranker baselines.

\begin{table}[h]
\centering
\caption{\textsc{ProcessBench} external-source summary. Higher is better for
AUROC and validation-selected test accuracy. Best values are bold where
applicable.}
\label{tab:appendix_processbench}
\small
\begin{tabular}{lcccc}
\toprule
model & task & AUROC & acc@0.5 & acc@val-thr \\
\midrule
\texttt{frozen\_visible} & PB-Trace & 0.7863 & 0.7059 & --- \\
\texttt{frozen\_masked} & PB-Trace & 0.8075 & 0.7157 & --- \\
\texttt{frozen\_step\_only} & PB-Trace & \textbf{0.8759} & \textbf{0.7784} & --- \\
\texttt{reranker8\_visible} & PB-Trace & 0.7337 & 0.3490 & 0.6980 \\
\texttt{reranker8\_masked} & PB-Trace & 0.7229 & 0.3529 & 0.7000 \\
\texttt{frozen\_visible} & PB-Prefix & 0.6759 & 0.6229 & --- \\
\texttt{frozen\_masked} & PB-Prefix & 0.6813 & 0.6386 & --- \\
\texttt{frozen\_step\_only} & PB-Prefix & \textbf{0.6918} & 0.6375 & --- \\
\bottomrule
\end{tabular}
\end{table}

\begin{table}[h]
\centering
\caption{\textsc{ProcessBench} external-source counterfactual checks. For
answer-only swap, lower mean absolute delta indicates less answer sensitivity.
For local repair, higher values indicate better local discrimination.}
\label{tab:appendix_processbench_cf}
\small
\begin{tabular}{lccc}
\toprule
check & visible & masked & note \\
\midrule
answer-only swap mean abs.\ delta & 0.3226 & 0.0645 & 31 answer-swapped pairs \\
local repair discrimination & 0.3750 & 0.5000 & 8 answer-matched repair pairs \\
\bottomrule
\end{tabular}
\end{table}

The external-source counterfactual checks matter because they separate two
questions. The answer-only swap set shows that answer sensitivity transfers
cleanly to externally sourced traces rather than being an artifact of the
generated quartets. The answer-matched local-repair subset now materializes
fully on a small selected subset, so its current negative-to-inconclusive read
can no longer be dismissed as mere construction failure. Instead, the present
external-source picture is sharper: visible answer sensitivity persists, but
external local repair still does not outperform the masked view. The package
artifacts above therefore serve two purposes at once: they make the
external-source benchmark reproducible, and they make its current claim
boundary explicit rather than leaving it implicit in scattered experiment
logs.

\subsection{Negative repair results that shape the claim boundary}

Several method lines remain informative negative results rather than positive
headline stories:
\begin{itemize}[leftmargin=1.5em]
\item \textbf{Hard-neg replay:} utility can improve, but answer sensitivity
often rises with it, leaving a clear frontier rather than a dominant method.
\item \textbf{Dual-head separation attempts:} the minimal dual-head and the
step-scan variant both fail to produce a stable additional channel at the
current scale.
\item \textbf{Earlier cleaner-benchmark hardening:} the older iterative
hardening line was useful for discovering artifact failure modes, but
\textsc{Craft-Clean} is now the active cleaner benchmark regime.
\end{itemize}

\subsection{External blind-audit status}

The remaining benchmark-validation blocker is external human blind audit on the
\textsc{Craft-Clean} packet. The packet, response form, hidden key, and scoring script
already exist locally; what remains is reviewer return rather than additional
local infrastructure.
